\chapter{Un ordinateur en cellules vivantes}
% version complète: chapters/21_living.tex

\pencilfig{21}{\pencilart{21}}{Des neurones vivants sur une puce: sous eux, des milliers d'électrodes, et chacune sait écouter comme parler.}

Peut-on construire une machine non pas en silicium, mais en neurones vivants? Oui, et on la construit déjà. On fait pousser des neurones directement sur une puce à milliers d'électrodes, en couche plane ou en petites boules qui ressemblent à un cerveau minuscule. Chaque électrode sait à la fois écouter et envoyer un signal. C'est un banc d'essai à schéma ouvert: tout peut être mesuré, et on peut intervenir partout. Seulement, il n'y a pas de stations de radio: c'est un bus de données sans registres de contrôle.

\section*{Ce que le réseau fait tout seul}

Livrée à elle-même, la culture s'embrase de temps en temps tout entière: presque tous les neurones cliquent ensemble, puis se calment, à des intervalles d'une seconde à plusieurs minutes. C'est l'oscillateur que nous connaissons déjà: le réseau s'excite lui-même et se fatigue, jusqu'à ce que les connexions reconstituent leur réserve de substance.

\pencilfig{129}{%
% spike raster of 8 neurons in a culture left to itself: sparse single clicks and shared bursts
\foreach \b in {0.9,3.1,4.0,6.6}{\fill[pcAmber, opacity=0.18] (\b-0.1,-0.15) rectangle (\b+0.32,2.95);}
\foreach \y/\ss in {0/{0.96,1.0,1.13,2.43,3.0,3.12,3.24,4.02,4.09,4.15,6.66,6.69,6.81},
  1/{0.46,0.88,1.0,1.05,3.09,3.2,3.94,4.04,4.37,6.65,6.71},
  2/{0.73,0.84,0.94,1.41,3.08,3.12,3.21,4.05,4.06,4.17,6.59,6.63},
  3/{0.89,0.93,1.06,3.1,3.19,3.28,3.89,3.94,4.13,4.2,4.31,6.63,6.65,6.78},
  4/{0.87,0.96,3.09,3.15,3.23,3.73,3.96,4.06,4.19,5.98,6.66,6.68,6.75},
  5/{0.44,0.92,0.97,3.05,3.22,3.27,4.01,4.07,4.17,5.76,6.57,6.73,6.75},
  6/{0.92,1.0,1.73,2.85,3.18,3.19,4.0,4.07,6.53,6.66,6.73},
  7/{0.87,0.92,3.13,3.13,3.27,3.99,4.06,4.17,6.44,6.61,6.72,7.13}}{
  \foreach \s in \ss {\draw[pcBlue, line width=0.7pt] (\s,0.35*\y) -- (\s,0.35*\y+0.25);}}
\draw[arr] (0,-0.3) -- (7.5,-0.3); \node[lbl, anchor=north] at (3.75,-0.32) {temps};
\node[lbl, anchor=east, align=right] at (-0.05,1.35) {huit\\neurones};
\node[lbl, text=pcAmber!70!black, anchor=south] at (3.55,2.95) {salves communes};
}{Une culture sur puce livrée à elle-même. Les clics isolés sont rares, et de temps en temps presque tout le réseau s'embrase d'un coup, puis se calme jusqu'à ce que les connexions aient reconstitué leur réserve.}

\section*{Un professeur sans dopamine}

Comment instruire un tel réseau s'il n'y a pas de dopamine? Les premières méthodes étaient électriques. Par exemple: stimuler le réseau jusqu'à ce qu'il donne la réponse voulue, puis s'arrêter aussitôt. La réponse se fixe. Dans une autre expérience, une culture pilotait un \og corps\fg{} virtuel simple et apprenait à le mener vers un but. Dans toutes ces expériences, le professeur est un motif de courant choisi par l'ingénieur.

\section*{De la dopamine par un éclair de lumière}

Un jour, on a essayé de rendre le professeur chimique. Sur l'une des plateformes à petites boules de neurones vivants, on a ajouté au liquide nutritif de la dopamine enfermée dans une \og cage\fg{} photosensible qui l'empêche d'agir. Quand le réseau répondait à un stimulus plus fort qu'avant, on le \og récompensait\fg{}: un éclair d'ultraviolet brisait la cage, et la dopamine était libérée. Au fil de deux cent quarante répétitions, les réponses ont augmenté dans l'ensemble. Mais les auteurs eux-mêmes écrivent honnêtement qu'une seule expérience de ce genre ne permet pas de dire si c'est bien la dopamine qui en est la cause.

\section*{Le pendule sur un chariot}

Dans une expérience récente, on a appris à de petites boules de neurones à maintenir un pendule en équilibre sur un chariot, un problème classique pour les programmes d'apprentissage. Le choix des impulsions d'entraînement revenait à un programme. Les signaux choisis par le programme marchaient mieux que des signaux aléatoires, mais après trois quarts d'heure de repos, l'acquis disparaissait. Et si l'on bloquait les connexions \nt{GLUT}, l'apprentissage ne se faisait pas du tout: c'est bien là que loge l'acquis. Le professeur reste pourtant à l'extérieur; seulement, ce n'est plus un ingénieur, mais un programme.

\section*{Le réservoir vivant}

Il existe une autre approche: ne pas instruire du tout le réseau vivant. On s'en sert comme d'un \og réservoir\fg{}, un système riche et enchevêtré dont un simple compteur électronique apprend à déchiffrer les réponses. C'est ainsi qu'une petite boule vivante a aidé à distinguer des voix. Cela dit, selon les auteurs, les connexions à l'intérieur de la boule changeaient elles aussi pendant le travail.

\section*{Le coût et la question}

Un million de neurones sur une puce consomment environ un dix-millième de watt, bien moins que l'électronique qui les entoure. Et le banc d'essai ne sait pas décider seul de ce qui compte comme un succès: ce rôle revient toujours à quelqu'un d'extérieur. À mesure que ces systèmes grandissent se pose aussi une question éthique, que personne n'a encore résolue.

\fullversion{21}
